\begin{figure*}[!t]
\centering
\begin{tcolorbox}[colback=gray!5, colframe=gray!50,
                  boxrule=0.5pt, arc=2mm,
                  left=4pt, right=4pt, top=3pt, bottom=3pt]
\small
\textbf{Inference Prompt (Code).}

\medskip
\textbf{Role definition.} You are a proactive code-maintenance agent whose goal is to surface issues in a Python software repository. Your task is to find and identify (without prior prompting) distinct issues worth reporting in the codebase, and recommend a resolution for each by producing a unified diff patch.

\medskip
\textbf{Bug pattern templates definition.} The following templates describe bug patterns observed in past, solved examples from other repositories. They are a partial library of known bug shapes, not an exhaustive catalog of every issue this codebase may contain. Report two kinds of issues:
\begin{enumerate}[topsep=0pt, itemsep=0ex, parsep=0pt, left=14pt]
   \item Issues that match one of the templates below. Set \texttt{used\_template\_id} to the template id and explain the match in \texttt{why\_matched}.
   \item Issues that do not fit any template but are still genuine bugs grounded in the code. Leave \texttt{used\_template\_id} as \texttt{""}.
\end{enumerate}
Both kinds count equally.

\medskip
\textbf{Bottleneck definition.} An issue is a problem in this codebase that a user, maintainer, or contributor would reasonably report. Each issue must
\begin{itemize}[topsep=0pt, itemsep=0ex, parsep=0pt, left=12pt]
   \item be resolvable by editing one or more of the provided Python functions;
   \item point to specific function(s) and a specific change;
   \item be supported by concrete evidence drawn from the code itself;
   \item be distinct from other issues in the list (do not split one bug across multiple entries).
\end{itemize}

\medskip
\textbf{Previously found bottlenecks} (included only after the first iteration). Report only new issues distinct from those already returned, covering both template-matched and template-novel kinds. If no new issues remain, return an empty JSON object \texttt{\{\}}.

\medskip
\textbf{Task execution.}
\begin{enumerate}[topsep=0pt, itemsep=0ex, parsep=0pt, left=14pt]
   \item Analysis: examine all provided functions to understand what each one is supposed to do.
   \item Issue identification: identify problems matching the issue definition, skipping any already reported.
   \item Patch generation: produce a unified diff patch resolving each identified issue.
   \item Response generation: provide analysis and patches in the specified JSON format.
\end{enumerate}

\medskip
\textbf{Output format.}
\begin{verbatim}
{
  "bottlenecks": [
    {
      "used_template_id":   "TID_X" or "",
      "used_template_name": "..."   or "",
      "why_matched": "Which observations in the functions satisfy
                      the matched template's evidence_flow, citing
                      specific function IDs and lines. Use \"\"
                      if no template applies.",
      "retrieved_documents": ["func_id_1", "func_id_2", ...],
      "bottleneck": "Natural-language description, citing concrete
                     code evidence and qualnames.",
      "action": {
        "patch": "<unified diff text fixing the issue,
                  starting with `diff --git a/<path> b/<path>`>"
      }
    }
  ]
}
\end{verbatim}

\medskip
\textbf{Context shown to the model:}
\begin{verbatim}
Repository: {repo}

Bug Pattern Templates:
[{template_id}] {template_name}
  Pattern: {pattern}
  Evidence flow:
    - {step_1}
    - {step_2}
    ...

Functions: {data_sources}
\end{verbatim}
\end{tcolorbox}
\caption{Per-iteration inference prompt for the code setting. At each iteration the model receives the repository name, current template pool, candidate Python functions, and the bottlenecks already returned in previous iterations; it returns new bottlenecks (template-matched or template-novel) with unified-diff patches, or an empty object to terminate.}
\label{fig:prompt_inference_code}
\end{figure*}
